\section{Thiết kế RTOS}

Phiên bản firmware hiện tại chạy tuần tự trong Arduino \texttt{loop()}. Khi chuyển sang RTOS, các chức năng nên tách thành task để đảm bảo timing rõ ràng và tránh một thao tác I/O làm chậm toàn bộ hệ thống.

\begin{table}[H]
    \centering
    \caption{Đề xuất phân chia task RTOS}
    \begin{tabularx}{\textwidth}{>{\bfseries}l X X}
        \toprule
        Task & Chức năng & Dữ liệu trao đổi \\
        \midrule
        \texttt{Task\_Sensor} & Đọc PMS7003, Modbus, VOC theo chu kỳ & Gửi \texttt{SensorSample} vào queue/store \\
        \texttt{Task\_Process} & Đánh giá IAQ, hysteresis, hold-time & Nhận \texttt{SensorSample}, xuất \texttt{IaqState} \\
        \texttt{Task\_Control} & Cập nhật LED/thiết bị chấp hành & Nhận \texttt{IaqState} \\
        \texttt{Task\_MQTT} & Publish telemetry và nhận lệnh điều khiển & MQTT payload JSON \\
        \texttt{Task\_Log} & Ghi CSV cục bộ hoặc dump log & \texttt{SensorSample} theo chu kỳ \\
        \bottomrule
    \end{tabularx}
\end{table}

\texttt{SnapshotStore} trong source hiện tại đã dùng critical section để bảo vệ mẫu đo mới nhất. Đây là nền tảng phù hợp khi nhiều task cùng đọc/ghi dữ liệu.

\section{Thiết kế MQTT}

Theo contract chung, hệ thống dùng các topic:

\begin{itemize}[leftmargin=2em]
    \item \texttt{iot/sensor/data}: ESP32 publish telemetry lên backend.
    \item \texttt{iot/device/control}: backend publish lệnh điều khiển xuống ESP32.
\end{itemize}

Payload dữ liệu cần chứa \texttt{device\_id}, timestamp, dữ liệu cảm biến và cờ trạng thái. Với firmware hiện tại, payload nên mở rộng từ \texttt{SensorSample} để có đủ PM, CO2, VOC, nhiệt độ, độ ẩm và cờ hợp lệ thay vì chỉ gửi temperature/humidity.

\section{Cơ chế vận hành khi mất kết nối}

ESP32 không nên phụ thuộc hoàn toàn vào server. Khi MQTT hoặc Wi-Fi mất kết nối, thiết bị vẫn đọc cảm biến, đánh giá IAQ, điều khiển cục bộ và ghi log SPIFFS. Khi kết nối khôi phục, hệ thống có thể tiếp tục publish dữ liệu realtime; phần đồng bộ log lịch sử là hướng mở rộng.
